% =============================================================================
% Partie 0 — Introduction
% =============================================================================
\section{Introduction}

\subsection{Objet du document}

Ce dossier décrit la conception d'un éditeur de documents partagé déployé sur
AWS. L'application poursuit deux buts que l'énoncé oppose volontairement :

\begin{itemize}[leftmargin=1.4em,itemsep=2pt]
  \item plusieurs centaines de personnes écrivent \emph{en même temps} dans le
        même texte, et aucune frappe acceptée ne doit être perdue ;
  \item un lien public distribue le document en PDF à \emph{exactement} $N$
        personnes, jamais $N+1$, chacune recevant le document tel qu'il était
        au moment où elle a réservé son exemplaire.
\end{itemize}

Le premier objectif demande de la souplesse : des écritures concurrentes, un
historique, de la reprise après coupure. Le second demande de la rigidité : un
compteur exact, une transaction, une réponse unique. Tenir les deux dans la même
application est l'essentiel du travail de conception.

\subsection{Ce que ce dossier est, et n'est pas}

Ce dossier décrit la conception \emph{réellement implémentée}, vérifiée contre
le code source, et non une intention de conception. Trois conséquences :

\begin{enumerate}[leftmargin=1.4em,itemsep=2pt]
  \item là où l'implémentation s'écarte du plan initial, l'écart est signalé avec
        sa raison — ces écarts sont souvent la partie la plus instructive ;
  \item chaque affirmation sur une contrainte renvoie à un test unitaire, un
        script d'intégration ou un scénario de charge, tous exécutables ;
  \item les limites sont énoncées sans être atténuées. Une limite annoncée vaut
        mieux qu'une limite découverte par le lecteur.
\end{enumerate}

\subsection{Plan}

La \textbf{partie A} traite de l'architecture : contexte, déploiement, décisions
et alternatives écartées, flux internes, sécurité, infrastructure, limites.

La \textbf{partie B} traite de la conception applicative : cas d'utilisation,
organisation des paquets, modèle du domaine, états, séquences système, puis une
lecture par contrainte, de C1 à C6.

La \textbf{partie C} conclut : ce qui est prouvé, ce qui reste à mesurer, et ce
que ce projet a réellement appris.

% -----------------------------------------------------------------------------
\section{Contexte, objectif et contraintes}

\subsection{Contexte}

Le sujet impose de démontrer la maîtrise des services AWS au travers d'une
application non triviale, dont chaque choix doit pouvoir être justifié. Il
précise que « une application simple, mais dont chaque choix est compris et
justifié, vaut mieux qu'une application complexe que vous ne savez pas
expliquer ». Ce dossier suit cette consigne : peu de services, mais chacun
employé pour ce qu'il fait de mieux, et chaque fois avec l'alternative écartée.

Deux points du sujet sont éliminatoires : les contraintes C1 et C2. Une
application rapide qui les viole est non fonctionnelle, quel que soit son
polissage.

\subsection{Les six contraintes}

\begin{table}[H]
\centering
\small
\begin{tabular}{@{}p{0.6cm}p{5.3cm}p{7.3cm}@{}}
\toprule
\textbf{\#} & \textbf{Contrainte de l'énoncé} & \textbf{Réponse retenue} \\
\midrule
C1 & Exactement $N$ téléchargements, jamais $N+1$
   & transaction DynamoDB : décrément de \code{remaining} et création du ticket
     dans la même écriture conditionnelle. Une garde en mémoire allège la base
     sous le pic, mais n'est jamais l'autorité \\
\addlinespace
C2 & Aucune modification acceptée perdue ; convergence
   & CRDT Yjs avec \emph{un seul écrivain} par document ; une modification n'est
     acquittée qu'\emph{après} écriture durable dans le journal \\
\addlinespace
C3 & Document tel qu'il est à $T_1$, complet et cohérent
   & le claim fige un numéro de séquence $T_1$ ; le PDF est construit depuis le
     journal jusqu'à ce numéro, une seule fois par version, de façon asynchrone \\
\addlinespace
C4 & Jamais plus que son crédit, même sur plusieurs appareils
   & compteur DynamoDB par (e-mail, jour) débité dans la même transaction que
     l'écriture au journal ; comptage effectué côté serveur, sur le binaire Yjs \\
\addlinespace
C5 & 15\,000 clics (80\,\% en 60\,s) pendant que 300 rédacteurs écrivent
   & micro-cache nginx sur l'état du lien, garde en mémoire sur les
     réservations, diffusion par lots de 250\,ms, instance dimensionnée par la
     mesure. Trois états explicites : prêt, patienter, épuisé \\
\addlinespace
C6 & Dernier exemplaire : visiteurs prévenus vite, plus de travail inutile
   & état épuisé monotone : réponse immédiate, aucun appel à la base, mis en
     cache longuement par nginx \\
\bottomrule
\end{tabular}
\caption{Les six contraintes et la réponse architecturale retenue.}
\label{tab:contraintes}
\end{table}

\subsection{Hypothèses d'interprétation}

Le sujet laisse plusieurs points ouverts. Les choix ci-dessous ne sont pas
confirmés par l'enseignant : ils sont donc explicités, pour que le lecteur
puisse contester l'interprétation plutôt que le résultat.

\begin{table}[H]
\centering
\small
\begin{tabular}{@{}p{5.4cm}p{10.2cm}@{}}
\toprule
\textbf{Sujet} & \textbf{Interprétation retenue} \\
\midrule
Consommation d'un exemplaire & au clic qui réserve le téléchargement, et non au
  téléchargement effectif du fichier \\
\addlinespace
Instant $T_1$ & lecture du dernier numéro de séquence au début du claim \\
\addlinespace
Crédit & 10\,000 caractères par jour ; insertions \emph{et} suppressions
  comptées ; restauration de version gratuite \\
\addlinespace
Recharge du crédit & minuit, heure de Douala — et non minuit UTC (voir B.6) \\
\addlinespace
Blocage & 4 refus pour crédit épuisé dans la journée bloquent l'écriture
  jusqu'à la recharge ; la lecture et le téléchargement restent possibles \\
\addlinespace
Pic d'écriture & 300 rédacteurs dans le \emph{même} document : c'est le pire cas,
  retenu comme cible ; le palier réellement atteint est celui qui est rapporté \\
\addlinespace
Identité & nom et e-mail, sans mot de passe. Le sujet l'autorise ; l'e-mail est
  la clé du crédit et du partage \\
\addlinespace
Livrables & français, dépôt GitHub privé, jury ajouté en collaborateur \\
\bottomrule
\end{tabular}
\caption{Hypothèses d'interprétation, à confronter à l'énoncé.}
\label{tab:hypotheses}
\end{table}

\hypothese{La recharge du crédit à minuit \emph{heure de Douala} s'écarte de
l'hypothèse initiale (minuit UTC). L'écart vient de l'usage : un utilisateur à
Douala lisait « recharge à 01:00 » et ne comprenait pas. Le décalage est
constant toute l'année, le Cameroun n'appliquant pas d'heure d'été, ce qui rend
l'arithmétique exacte sans table de fuseaux.}

\subsection{Objectif mesurable}

La conception n'est acceptable que si elle est mesurable. Quatre jalons ont
guidé les décisions :

\begin{enumerate}[leftmargin=1.4em,itemsep=2pt]
  \item $N$ réservations exactement pour un lien de $N$ exemplaires, sous
        15\,000 clics dont 80\,\% en 60 secondes, sans aucune erreur serveur ;
  \item zéro modification acceptée sans acquittement, sous écriture simultanée
        de plusieurs dizaines de rédacteurs ;
  \item convergence du contenu vérifiable : journal sans trou, reconstruction
        identique par deux chemins indépendants ;
  \item un palier de rédacteurs réellement atteint, rapporté tel quel, même s'il
        est inférieur à la cible.
\end{enumerate}

Le troisième point mérite d'être souligné : un test de charge mesure un
\emph{débit d'acquittement}, pas une convergence de contenu. Les deux preuves
sont distinctes, et le dossier les traite séparément.
